\chapter{Particles, Alternatives, and Scope}\label{ch:particles}\label{ch:function-words}

\section{Particle taxonomy and structural positions}\label{sec:adv-particles}

\label{sec:adv-particles-general}

Adverbial particles form a closed series whose members differ in scope,
position, and prosodic dependence. The phasal particles \ird{vy} and \ird{so}
are unstressed proclitics supported by the following predicate and remain
outside its lexical stress domain, so their syllables do not enter the
calculation of subsidiary footing (§\,\ref{sec:secondary-stress}).
Comment-domain particles may be weakly attached in connected speech, whereas
stance particles can form independent prosodic words with ordinary initial
stress. Deaccentuation in a phrase does not by itself make an independent
particle a clitic (§\,\ref{sec:stress}).

The semantic classes below do not correspond one for one to syntactic heights.
Reportative \ird{či} occupies an outer evidential position, whereas
discourse-reorienting \ird{jas} follows additive \ird{oh} and precedes
phasal \ird{vy}/\ird{so}. In ordinary finite clauses, overt arguments
and lexical modifiers precede the entire articulated particle field, including
stance particles and interrogative \ird{by}; their lower chain occurrences
retain the relations described in §\,\ref{sec:th-particle-framework}.

The traditional Iridian term for these items is \ird{zemo} `grain, particle'.
The label \term{adverbial particle} is retained here because many of them behave
broadly like clause adverbs, but the class is wider than ordinary manner or
temporal adverbs. Particles may contribute phasal contour, addition,
affirmation, reportative force, evidential stance, mirativity, and other
clause-level meanings. What unifies them is not a single semantic type but a
distribution across ordered structural zones, with independently licensed
interpretations within those zones.

One element that appears in a similar linear region but belongs to a different
syntactic class is the interrogative particle \ird{by}. It is treated in
§\,\ref{sec:by-pol-head}, where its status as the overt exponent of \synhead{Pol\textsuperscript{0}} —
a dedicated clause-typing head above the comment domain — is developed. It
is separated from the adverbial particle classes here because grouping it with
them would misrepresent its structural position.

\begin{table}[ht]
	\footnotesize
	\caption{Adverbial particle classes.}\label{tab:adv-particle-classes}
	\begin{tabularx}{\textwidth}{l L L l}
		\toprule
		{\sc class} & {\sc members} & {\sc core function and scope} & {\sc position} \\
		\midrule
		Class 1 &
		\ird{vy}, \ird{so} &
		phasal or scalar comparison over a polarised predication or licensed alternatives &
		closest to predicate \\
		Class 2 &
		\ird{oh}, \ird{mó}, \ird{či}, \ird{kamo}, \ird{dlí}, \ird{jas} &
		propositional and illocutionary meanings such as addition, affirmation, hearsay stance, exclusivity, temporary priority, and contrast &
		mixed (see text) \\
		Class 3 &
		\ird{oče}, \ird{nadě}, \ird{hladě}, \ird{kady}, \ird{mlada} &
		speaker stance over the whole utterance, including mirative, consequential, inferential, speculative, and hypothetical meanings &
		farthest from predicate \\
		\bottomrule
	\end{tabularx}
\end{table}

The unmarked order, with any unexpressed positions omitted, is:

\pex
\a \ird{external frame / marked information structure > ordinary arguments and modifiers > particle field > predicate}
\xe

\noindent Within the particle field, the established partial order is
outer speech-act/evidential material $>$ \ird{by} $>$ \ird{oh} $>$
\ird{jas} $>$ \ird{vy}/\ird{so}, with absent positions omitted.
This order does not license every combination or settle every other operator's
relative position. Ordinary placement is distinct from aboutness topicalisation:
a quantificational pivot can precede particles without becoming a topic.
Pol\textsuperscript{0} remains locally head-initial, although ordinary
constituents are pronounced above it. Particle specifiers and silent final
heads retain their independently defined directionality. Phonological
attachment follows the adjacent surface host; the contiguous field need not
form a single phonological word, and independently pronounced stance particles
retain initial lexical stress. Nominal predicates receive a separate treatment,
including the negative equative order \ird{Marek doktor vy staly}
(§\,\ref{sec:th-particle-framework}).

Derived manner adverbs in \ird{-ě}, such as \ird{skorně} \trsl{quickly},
occur before the particle field through ordinary placement, while their lower
modifier positions lack the fixed functional ordering of the closed series.

\pex
\a
\begingl
\gla Janek vy pjášček.//
\glb Janek.\Dir{} already eat-\Antip{}-\Pf{}//
\glft \trsl{Janek has already eaten.}//
\endgl
\a
\begingl
\gla Tak tóm Janku oče by vy koupnik?//
\glb this book.\Dir{} Janek.\Trs{} mirative question already buy-\Trs{}-\Pf{}//
\glft \trsl{Has Janek, surprisingly, already bought this book?}//
\endgl
\xe

\section{Phasal \ird{vy} and \ird{so}}\label{sec:class1-particles}

The predicate-adjacent particle position admits either \ird{vy} or \ird{so},
never both. These operators relate a proposition to a contextually ordered
domain of stages or alternatives; they do not realise lexical aspect or the
viewpoint aspect expressed by Asp.Align. Their interpretation requires a recoverable
comparison domain. A reference time supplies such a domain in ordinary phasal
uses, while an explicitly established series, threshold, or set of alternatives
licenses the scalar and additive uses described below. Neither particle
licenses an arbitrary ordering merely because an English translation permits it.
Persistence also has an established nonfinite use with the state-frame converb
in \ird{so byleň} \trsl{while still a child} (§\,\ref{sec:state-frame}).

The distinction between ordinary phasal polarity and scalar addition has
comparative precedents, although their distribution is language-specific.
Persohn distinguishes ordinary preceding-stage polarity from informative
scalar alternatives in German, without making early attainment obligatory
\cite[§§2.1--2.2, 6]{persohn2024scalar}. The temporal approximation and the
separate licensing conditions adopted here make that same analytical distinction;
they do not identify the Iridian particles with German lexical entries.

\subsection{The common opposition and negative predication.}
In the ordinary temporal approximation, let $P$ be the polarised predication,
$r$ the reference stage, and $r^{-}$ its contextually relevant preceding stage.
Affirmative \ird{vy} asserts $P(r)$ against preceding $\neg P(r^{-})$;
\ird{so} asserts $P(r)$ against preceding $P(r^{-})$. The preceding stage
belongs to the salient comparison, rather than to any arbitrarily remote
period. The preceding-stage condition supplies presupposed phasal background in
these ordinary uses, distinct from the asserted reference-stage proposition.
These formulae characterise ordinary attainment and persistence;
they are not complete lexical entries deriving every scalar use.

Negation belongs inside this phasal relation: \ird{vy}(\neg P) presents the
negative state as attained after positive $P$, whereas \ird{so}(\neg P)
presents that negative state as continuing. Thus \trsl{no longer} and
\trsl{not yet} follow from the same opposition as \trsl{already} and
\trsl{still}. An expectation of later attainment under \trsl{not yet} is
defeasible: continued non-attainment does not itself assert eventual success.
Nor does \ird{vy} entail that attainment was earlier than expected.

\begin{table}[ht]
\caption{The ordinary phasal opposition.}\label{tab:phasal-polarity}
\begin{tabularx}{\textwidth}{l L L}
\toprule
Construction & Relevant preceding stage & Reference stage \\
\midrule
\ird{vy} $P$ & $\neg P$ & $P$: already attained \\
\ird{so} $P$ & $P$ & $P$: still continuing \\
\ird{vy} $(\neg P)$ & $P$ & $\neg P$: no longer \\
\ird{so} $(\neg P)$ & $\neg P$ & $\neg P$: not yet \\
\bottomrule
\end{tabularx}
\end{table}

\begin{figure}[ht]
\centering
\begin{tikzpicture}[
  x=1cm,y=1cm,
  every node/.style={font=\small},
  phaslabel/.style={anchor=east},
  truth/.style={draw,rounded corners,minimum width=1.3cm,minimum height=.55cm}
]
% Arguments: row y, particle, scope, preceding P-value, reference P-value, reading.
% P is the positive base proposition in every row, including the negative rows.
\newcommand{\PhasalRow}[6]{%
  \node[phaslabel] at (0,#1) {\ird{#2} $#3$};
  \node[truth] (before#2#1) at (1.4,#1) {$#4$};
  \node[truth] (now#2#1) at (4.0,#1) {$#5$};
  \draw[->] (before#2#1) -- (now#2#1);
  \node[anchor=west] at (5.0,#1) {#6};
}
\node at (1.4,.7) {preceding stage};
\node at (4.0,.7) {reference stage};
\PhasalRow{0}{vy}{P}{\neg P}{P}{already}
\PhasalRow{-1}{so}{P}{P}{P}{still}
\PhasalRow{-2}{vy}{(\neg P)}{P}{\neg P}{no longer}
\PhasalRow{-3}{so}{(\neg P)}{\neg P}{\neg P}{not yet}
\end{tikzpicture}
\caption{Ordinary phasal polarity. The boxes give the truth value of the
positive base proposition $P$, while the left-hand labels show the proposition
within the particle's scope. The preceding stage is the contextually relevant
predecessor, not an arbitrarily remote occasion. The reference-stage proposition
is asserted; the preceding-stage contrast or persistence supplies its phasal
background. These time-based schemata do not by themselves derive scalar,
additive, or projected interpretations. In particular, \trsl{not yet} does not
assert a later transition to $P$.}
\label{fig:phasal-polarity}
\end{figure}


The physical-state schemata below use the lexical-slot convention of
§\,\ref{sec:lexical-slot-schemata}. Each predicate slot denotes a complete
finite form of the indicated aspect and polarity, rather than an uninflected
English stem; no declension class or stative stem formation is assigned to
the missing lexeme.
\pex
\a
\begingl
\gla \textnormal{[[western-gate]]}\textsubscript{\Dir{}} vy \textnormal{[[be-open]]}\textsubscript{\Cont{}}.//
\glb western.gate.\Dir{} already be.open.\Cont{}//
\glft \trsl{The western gate is already open.}
\lingcontext{It was closed at the preceding stage.}//
\endgl
\a
\begingl
\gla \textnormal{[[oldest-candle]]}\textsubscript{\Dir{}} so \textnormal{[[burn]]}\textsubscript{\Prog{}}.//
\glb oldest.candle.\Dir{} still burn.\Prog{}//
\glft \trsl{The oldest candle is still burning.}
\lingcontext{It was burning at the preceding stage.}//
\endgl
\a
\begingl
\gla \textnormal{[[eastern-road]]}\textsubscript{\Dir{}} vy \textnormal{[[be-flooded]]}\textsubscript{\textsc{neg.cont}}.//
\glb eastern.road.\Dir{} already be.flooded.\Neg{}.\Cont{}//
\glft \trsl{The eastern road is no longer flooded.}
\lingcontext{It was flooded at the preceding stage.}//
\endgl
\a
\begingl
\gla \textnormal{[[eastern-road]]}\textsubscript{\Dir{}} so \textnormal{[[be-flooded]]}\textsubscript{\textsc{neg.cont}}.//
\glb eastern.road.\Dir{} still be.flooded.\Neg{}.\Cont{}//
\glft \trsl{The eastern road is not flooded yet.}
\lingcontext{Its non-flooded state continues; later flooding is not entailed.}//
\endgl
\a
\begingl
\gla Pivo vy zada.//
\glb beer.\Dir{} already \NegExst{}-\Cont{}//
\glft \trsl{There is no more beer.} \lingcontext{Beer was available before.}//
\endgl
\a
\begingl
\gla Pivo so zada.//
\glb beer.\Dir{} still \NegExst{}-\Cont{}//
\glft \trsl{There is no beer yet.} \lingcontext{Its absence continues; arrival is not guaranteed.}//
\endgl
\xe

\subsection{States, events, thresholds and alternatives.}
An open gate supplies an attained state, while a burning candle supplies
a state whose earlier validity licenses persistence. A sufficiency use of
\ird{vy} similarly requires a contextual criterion. If the stipulated need is
beer for two guests, \ird{Pivo vy jeca} \trsl{There is already enough beer}
can mark that requirement as met; without a salient quantity threshold, the
same string merely marks the attained availability of beer. Sufficiency is
therefore not a meaning freely added to every existential predication.

With completed events, \ird{so} may extend an established series by a further
occurrence. A prior arrival followed by another arrival on Saturday licenses
the following additive reading. A single arrival, with no salient series or
maintained alternative, does not license the translation solely because much
time has since elapsed. This addition can describe events which also satisfy
the renewed-occurrence conditions of \ird{jas vy}, although addition alone
does not specify that construction's intervening contrary stage
(§\,\ref{sec:class2-particles}).

\pex
\begingl
\glpreamble \lingcontext{Janek visited on Friday; Saturday supplied another arrival.}//
\gla Janek sobotě so ščének.//
\glb Janek.\Dir{} Saturday.\Ind{} additional arrive-\Pf{}//
\glft \trsl{Janek arrived on Saturday too.}//
\endgl
\xe

Scalar \trsl{even} requires more than addition: the context must order
alternatives and identify the added one as unlikely or marginal. Visiting the
store can be a surprising final addition to an already specified series of
errands; without that ordering, \ird{so} supplies addition, not surprise.
Affirmative \ird{mó} ratifies the proposition, while mirative \ird{oče}
presents it as counterexpectational; neither function is the scalar ordering
itself.

\pex
\begingl
\glpreamble \lingcontext{Several unlikely errands were accomplished; a store visit was considered least likely.}//
\gla Magazinam de byh so stójek!//
\glb store.\Ind{} into yesterday additional go-\Pf{}//
\glft \trsl{I even went to the store yesterday!}//
\endgl
\xe

Marginal category inclusion instead concerns the boundary at which a predicate
remains applicable. In the English schema \trsl{The northern window still
counts as open}, a very small opening can make classification questionable
without presupposing another open window. An independently ordered set of
windows instead licenses scalar-additive \trsl{Even the northern window is
open}. Without either licensing context, the ordinary interpretation is
temporal persistence.
The distinction between boundary marginality and focus-sensitive scalar
addition is also explicit in Persohn's analysis of German
\cite[§§4.1--4.2]{persohn2024scalar}.

In content questions, additive \ird{so} extends a contextually supplied set of
answers. This yields \trsl{what else}, \trsl{who else}, \trsl{where else},
or \trsl{which else} according to the wh-expression and its alternatives,
without altering the interrogative syntax or making \ird{so} a question marker.

\pex
\begingl
\glpreamble \lingcontext{Marek has eaten bread; the speaker asks for another item.}//
\gla Ježe Marku so pjaštník?//
\glb what Marek.\Trs{} additional eat-\Trs{}-\Pf{}//
\glft \trsl{What else did Marek eat?}//
\endgl
\xe

\subsection{Viewpoint, modality and clause type.}
Asp.Align specifies viewpoint and alignment before the phasal operator is
interpreted. Perfective morphology does not force \ird{vy}, nor does it
exclude additive \ird{so}; progressive morphology can combine with either
attainment of an ongoing state or persistence of that state. A retrospective
supplies its own prior-event perspective, to which \ird{vy} adds attainment
at the relevant reference stage. Lexical preverbs remain within the complete
stem: traversal in \ird{mněkař-} and restitution in \ird{rustoj-} are
lexical contributions, not competing phasal particles
(§\,\ref{sec:lexical-preverbs}). Both particles retain their forms before
vowels, /v/, and sibilants, without an expanded allomorph.

Contemplative predicates contribute a projected or intended state. In
\ird{Janek sobotě vy ščenáž}, the arrival projected for Saturday has become
settled or imminent in the relevant plan; the particle alone does not assert
that Janek will have completed his arrival by Saturday. In
\ird{Janek sobotě so ščenáž}, that projection remains in force. Actual
completion at a future reference time requires an independently specified
viewpoint and reference structure. Likewise, \ird{mlada so} maintains a
desired or hypothetical possibility, not a presupposed actual occurrence.

Under imperative or hortative force, attainment concerns the urged transition
into action, rather than an asserted completed event. \ird{Mó vy pjaštím!}
\trsl{Go on, eat already!} urges initiation; \ird{pjaštím} is the regular
imperative, whose independently long Mood vowel remains unchanged. Urgency or
impatience follows from context or reinforcement by \ird{mó}, not from a
lexical entailment of \ird{vy}. Negative Mood likewise supplies a negative
prejacent within the directive interpretation, without requiring an indicative
\synhead{NegP} or treating the prohibited event as actual.

\begin{table}[ht]
\footnotesize
\caption{Phasal interpretation inside polar questions.}\label{tab:v-s-questions}
\begin{tabularx}{\textwidth}{l L L}
\toprule
Viewpoint & \ird{vy} & \ird{so} \\
\midrule
Perfective & \ird{Janek by vy pjášček?}\newline Has Janek already eaten? &
\ird{Janek by so pjášček?}\newline Did Janek eat on an additional occasion? \\
Progressive & \ird{Janek by vy pjaštíme?}\newline Is Janek already eating? &
\ird{Janek by so pjaštíme?}\newline Is Janek still eating? \\
Contemplative & \ird{Janek by vy pjaštáž?}\newline Has Janek's intention to eat become settled? &
\ird{Janek by so pjaštáž?}\newline Is Janek still intending to eat? \\
\bottomrule
\end{tabularx}
\end{table}

The perfective additive question requires a prior occasion; its bare form does
not entail surprise or successful action despite difficulty. The interrogative
operator takes the phasal proposition as its scope: \ird{by} supplies clause
typing, while \ird{vy} and \ird{so} supply the contrast being questioned
(§\,\ref{sec:th-polar-questions}).

\subsection{Interaction with other particles.}
Comment-domain particles ordinarily take scope over the phasal proposition,
while outer stance particles qualify the clause from higher positions within
the same surface field. The
combinations, their contextual requirements, and the two bounded exclusive
readings are brought together in §\,\ref{sec:particle-combinations}; the
individual entries below retain the distinctive readings of \ird{oh},
\ird{jas}, and the stance particles. Their composition preserves the
comparison domains specified here, including the distinction between actual
persistence and the continued availability of a desired possibility.

\section{Comment-domain and reportative operators}\label{sec:class2-particles}

The propositional particles contribute addition, affirmation, hearsay stance,
exclusivity, temporary priority, or discourse reorientation. Their positions
within the particle field differ: reportative \ird{či} occupies an outer
evidential position, while \ird{jas} follows \ird{oh} and precedes
\ird{vy}/\ird{so} (§\,\ref{sec:th-particle-framework}). The semantic
grouping is consequently broader than a single syntactic slot. Interrogative
\ird{by} has its own clause-typing
position (§\,\ref{sec:by-pol-head}).

\subsection{Oh} \ird{Oh} is primarily an additive particle. It presents the
proposition as holding for an additional discourse referent, or as an
additional fact alongside one already under discussion. In positive clauses it
is usually translated as \trsl{also} or \trsl{too}; under negation it often
corresponds to \trsl{either} or \trsl{neither.}

\pex
\a
\begingl
\gla Janek oh zdrávime.//
\glb Janek.\Dir{} also sleep-\Prog{}//
\glft \trsl{Janek is also sleeping.}//
\endgl
\a
\begingl
\gla Matuš oh zaščének.//
\glb Matuš.\Dir{} also \Neg{}-arrive-\Pf{}//
\glft \trsl{Matuš has not arrived either.}//
\endgl
\xe

With bounded predicates, especially perfectives and retrospectives,
\ird{oh} may develop a scalar phasal reading such as \trsl{finally} or
\trsl{at last}. This is best treated as a contextual extension rather than as a
separate lexical meaning: the particle still adds the event to an already
active discourse sequence, but the relevant addition is now a long-awaited
attainment. The reading is especially natural with \ird{vy}. With \ird{so} and a
negative predicate, \ird{so} supplies continuing non-attainment while
\ird{oh} requires an additional relevant proposition or participant; prolonged
waiting alone does not supply the additive contribution.

\pex
\a
\begingl
\glpreamble \lingcontext{Other guests have arrived; after a long wait,
the speaker sees Janek at the door as well.}//
\gla Janek oh vy ščének.//
\glb Janek.\Dir{} also already arrive-\Pf{}//
\glft \trsl{Janek has finally arrived.}//
\endgl
\a
\begingl
\glpreamble \lingcontext{The speaker has been waiting for Janek for some
time; Matuš has not arrived either, and Janek is an additional non-arrival.}//
\gla Janek oh so zaščének.//
\glb Janek.\Dir{} also still \Neg{}-arrive-\Pf{}//
\glft \trsl{Janek still has not arrived.}//
\endgl
\xe

\subsection{Mó} \ird{Mó}, pronounced \ipagloss{[ˈmuː]}, is an affirmative
or confirmatory particle. It marks
the speaker's ratification of the proposition as true, often where that truth
has been doubted, denied, questioned, or otherwise made salient in the
preceding discourse. Its stable contribution is therefore confirmatory rather
than emotional. Depending on context, the same particle may sound reassuring,
insistent, triumphant, annoyed, or reproachful.

\pex
\begingl
\glpreamble \lingcontext{A tells B that Janek has arrived. B doubts this, then
sees Janek in person and reports back.}//
\gla Janek mó vy ščének.//
\glb Janek.\Dir{} affirm already arrive-\Pf{}//
\glft \trsl{Janek really has arrived.}//
\endgl
\xe

In imperatives and hortatives, \ird{mó} keeps this ratifying force, but now
the speaker presents the urged course of action as one that should by now be
carried out. The result is often urgency, coaxing, or emphatic insistence.
With \ird{vy}, the reading commonly approaches \trsl{come on, already}.

\pex
\begingl
\gla Mó vy pjaštím!//
\glb affirm already eat-\Imp{}//
\glft \trsl{Go on, eat already!}//
\endgl
\xe

\subsection{Či} \ird{Či} is a hearsay or attributed-information particle. It
marks the proposition as one the speaker is passing on from another source
rather than presenting as newly verified by direct personal knowledge. The
source may be a prior report, common talk, something just heard, or information
already circulating in the discourse. Its effect is often lighter than a fully
formal reportative: the speaker indicates that the proposition is being relayed,
but need not thereby dramatize distance or strong scepticism.

\pex
\a
\begingl
\glpreamble \lingcontext{The speaker has just heard that Janek has arrived and
passes the information on.}//
\gla Janek či vy ščének.//
\glb Janek.\Dir{} reportedly already arrive-\Pf{}//
\glft \trsl{They say Janek has already arrived.}//
\endgl
\a
\begingl
\glpreamble \lingcontext{The speaker is repeating what they have been told about
Matuš.}//
\gla Matuš či zaščének.//
\glb Matuš.\Dir{} reportedly \Neg{}-arrive-\Pf{}//
\glft \trsl{They say Matuš has not arrived.}//
\endgl
\xe

\noindent What \ird{či} adds is not simply uncertainty. A speaker may use it
whether or not they personally find the proposition plausible; the essential
point is that the information is being framed as relayed or attributed rather
than as a fresh firsthand assertion.

\noindent This function should be kept distinct from complement constructions
built with an overt speech or cognition predicate. Those constructions
syntactically subordinate one proposition under another, as in \trsl{X said
that P} or \trsl{X thinks that P}. \ird{Či}, by contrast, does not embed the
proposition at all: it marks the speaker's stance toward the main clause
itself, presenting that clause as hearsay or relayed information. Accordingly,
\ird{či} is not redundant with quotative framing. The two constructions do
genuinely different jobs and may co-occur when both meanings are wanted:
syntactic attribution to a speech frame on the one hand, and speaker-level
hearsay stance toward the clause on the other.

\subsection{Kamo} \ird{Kamo} is an exclusive--scalar particle whose
interpretation depends on a focused associate and a contextually salient
question under discussion. With individual alternatives, it restricts who or
what satisfies a predicate; with ordered degrees, accomplishments or abilities,
it establishes an upper bound translated as \trsl{only}, \trsl{just} or
\trsl{merely}. These uses share an alternative-sensitive semantics, although
the question and its ordering differ.

The following English schemata specify meanings and focus domains; they
stand in place of lexical Iridian examples and carry no Iridian segmentation
or morphological analysis.
\pex
\a \trsl{Only the northern window is open.}
\lingcontext{The alternatives identify the relevant windows; the northern
window is open and the others are closed.}
\a \trsl{The child can read only the simplest passages.}
\lingcontext{The alternatives rank reading abilities by difficulty. Several
easy passages may be readable, although no higher level is attained.}
\xe

Japanese \textit{sika} and Korean \textit{pakkey} yield exclusive
readings in constructions requiring overt negation, whereas Japanese
\textit{dake} and Korean \textit{man} permit affirmative predicates
\cite[§\,3.1, examples~14--15]{sells2022}. Iridian \ird{kamo} also
permits an affirmative arrival assertion:

\ex\label{ex:kamo-arrival-counterpart}%
\begingl
\gla Karel kamo ščének.//
\glb Karel.\Dir{} only arrive-\Pf{}//
\glft \trsl{Only Karel arrived.}
\lingcontext{Karel arrived on the relevant occasion; the other relevant travellers did not}//
\endgl
\xe

The arrival predicate is affirmative, while exclusivity excludes the
alternative travellers. Negative \ird{Zide zaščének} in
(\ref{ex:concord-comparative-arrival}c) instead denies any arrival through
clausal negation and its dependent indefinite. When \ird{kamo} combines
with an independently negative predicate, the negative prejacent and the
exclusive comparison retain their separate contributions
(§\,\ref{sec:th-exclusive-phasal}).

The analysis follows Coppock and Beaver's distinction between a presupposed
lower bound and an asserted upper bound
\cite[§2.3, pp.~392--395]{coppockbeaver2014}. Let $C$ be the salient
alternative propositions, $p$ the prejacent, and $\succeq_C$ the independently
recoverable ordering of strength. The two contributions are
\[
\begin{aligned}
\mathrm{MIN}_C(p)(w)
  &=\exists q\in C\,[q(w)\wedge q\succeq_C p],\\
\mathrm{MAX}_C(p)(w)
  &=\forall q\in C\,[q(w)\rightarrow p\succeq_C q].
\end{aligned}
\]
\ird{Kamo} presupposes MIN and asserts MAX. MIN does not invariably
presuppose the prejacent itself: an attained stronger alternative can satisfy
the lower bound, although MAX excludes it when it exceeds $p$. Under a partial
ordering, MAX also excludes true alternatives incomparable with $p$; it is
therefore stronger than the prohibition of strictly higher alternatives alone.

For individual focus, the ordering can be made explicit by closing the
answers under sums of individuals. If $q_X$ says that all members of $X$ have
the property, define $q_X\succeq_C q_Y$ when $X\supseteq Y$. With
$p=q_{\{a\}}$, MIN requires that $a$ have the property, while MAX excludes
every true answer containing another relevant individual. This recovers the
ordinary exhaustive interpretation without arbitrarily ranking the individuals.
For scalar focus, the ordering instead compares levels on a salient scale:
the child's ability reaches the simplest level but no higher relevant level.
The ranking must be supplied by the context; \ird{kamo} does not establish
an unrestricted hierarchy of professions, achievements or values. In particular,
modest attainment need not express regret or disappointment.

Additional English schemata distinguish the cardinal, spatial and evaluative
scales without supplying Iridian lexical forms:
\begin{table}[ht]
\footnotesize
\caption{Supplementary exclusive and scalar schemata.}
\label{tab:kamo-supplementary-schemata}
\begin{tabularx}{\textwidth}{L L}
\toprule
English schema & Alternatives and interpretation \\
\midrule
The messenger brought only one letter. & Cardinal alternatives; no additional letters. \\
The ship sailed only as far as the southern island. & Ordered endpoints along a path. \\
The repairs consisted merely of replacing the roof. & Contextually ranked accomplishments; comparatively modest extent. \\
\bottomrule
\end{tabularx}
\end{table}

Focus also determines which material varies. In the English schema
\trsl{The gardener only watered the roses}, the alternatives may vary the
plants watered or the action performed, with the remaining material held
constant. The pronounced particle position does not select between those
associates. An existing lexical example illustrates individual focus within
a polar question:
\pex
\begingl
\gla Janek Roubžu na by kamo mulša?//
\glb Janek.\Dir{} Roubže.\Ind{} in question only live-\Cont{}//
\glft \trsl{Is Janek the only one living in Roubže?}//
\endgl
\xe
The locative phrase precedes \ird{by} through ordinary placement, while
the exclusive comparison remains within the question's comment domain.
The interpretation of the exclusive comparison within interrogation does
not independently settle projection of MIN; the lower-bound condition must
be distinguished from operator scope
(§\,\ref{sec:th-exclusive-phasal}).

\subsection{Dlí} \ird{Dlí} is a particle of temporary priority or interim
course. It presents the predication as the action to be carried out first, for
the time being, or in advance of some later step that remains contextually
available. In English it is often rendered by \trsl{first}, \trsl{for now},
\trsl{for the meantime}, or \trsl{just for the moment}. The particle frames the event as an immediate or provisional move within a
larger sequence of actions or decisions, whose later stages remain open.

\pex
\a
\begingl
\glpreamble \lingcontext{The speaker explains the immediate plan before deciding
what to do next.}//
\gla Dá dlí zdrava.//
\glb I.\Dir{} for.now sleep-\Cont{}//
\glft \trsl{I'll sleep for now.}//
\endgl
\a
\begingl
\glpreamble \lingcontext{The speaker urges the addressee to do the eating step
before anything else.}//
\gla Dlí pjaštím!//
\glb first eat-\Imp{}//
\glft \trsl{Eat first.}//
\endgl
\xe

\ird{Dlí} belongs to the propositional layer because
it manages the place of the whole predication within an intended sequence or
temporary plan, rather than merely describing the internal phase of the event
itself.

\subsection{Jas} \ird{Jas} /jas/ is a discourse-reorienting contrast particle. It
often introduces an alternative or contrast, but its use need not be confined
to strong semantic opposition. It may also mark topic shift, alternation,
turn-taking, or a looser \trsl{as for X} framing in which attention is
redirected to another participant or proposition.

\pex
\begingl
\gla Janek zdrava. Marek jas znóhtime.//
\glb Janek.\Dir{} sleep-\Cont{} Marek.\Dir{} contrast study-\Prog{}//
\glft \trsl{Janek is sleeping. Marek, for his part, is studying.}//
\endgl
\xe

This broader reorienting value also underlies corrective and defensive uses.
Here the speaker contrasts the current proposition not necessarily with an
explicitly stated opposite, but with some assumption, expectation, or implied
line of reasoning in the discourse.

\pex
\begingl
\glpreamble \lingcontext{The speaker's doctor tells them that they should start
eating more vegetables. They respond defensively.}//
\gla Dá graže jas pjašta.//
\glb I.\Dir{} vegetable.\Ind{} contrast eat-\Cont{}//
\glft \trsl{But I do eat vegetables.}//
\endgl
\xe

Because \ird{jas} reorients discourse rather than encoding contradiction by
itself, identical predicates are not automatically excluded. Such examples are
usually odd if no shift in participant set or discourse management is intended,
but they become natural when the second clause is interpreted as parallel,
alternating, or explicitly additive. Combination with \ird{oh} is especially
natural in this use.

\pex
\begingl
\gla Marek zdrávime. Janek oh jas zdrávime.//
\glb Marek.\Dir{} sleep-\Prog{} Janek.\Dir{} also contrast sleep-\Prog{}//
\glft \trsl{Marek is sleeping. Janek, meanwhile, is sleeping too.}//
\endgl
\xe

With \ird{vy}, \ird{jas} has a conventionalised renewed-occurrence reading,
translated as \trsl{again} or \trsl{once again}, which requires an earlier
occurrence and a relevant discontinuity before the current attainment.
The two particles retain distinct positions within the contiguous field:
ordinary arguments and modifiers precede both, while \ird{jas} precedes
predicate-adjacent \ird{vy}. In the departure example below, an earlier
departure is followed by a return and another departure; in the negative
example, an earlier failure to arrive is followed by a successful visit and
another failure. Negation is part of the recurring predication, so the
second example describes renewed non-arrival rather than uninterrupted
waiting for a first arrival.

\pex
\a
\begingl
\glpreamble \lingcontext{Janek was expected to stay for a while, but he has
departed before, returned, and now left once more.}//
\gla Janek jas vy váržek.//
\glb Janek.\Dir{} contrast already leave-\Pf{}//
\glft \trsl{Janek has left again.}//
\endgl
\a
\begingl
\glpreamble \lingcontext{Janek failed to arrive on an earlier occasion, came successfully on an
intervening visit, and has now failed to arrive again.}//
\gla Janek jas vy zaščének.//
\glb Janek.\Dir{} contrast already \Neg{}-arrive-\Pf{}//
\glft \trsl{Janek has once again not arrived.}//
\endgl
\xe

The renewed-occurrence construction is analysed here as adding a conventional
background condition to ordinary phasal attainment. Let $Q$ be the polarised
predication, $r$ its reference stage, and $r^{-}$ the relevant preceding
stage. Its declarative interpretation has the following profile:
\[
\begin{aligned}
\text{asserted content:}\quad & Q(r),\\
\text{background:}\quad & \neg Q(r^{-})\ \wedge\
   \exists r_0<r^{-}\,[Q(r_0)],\qquad r^{-}<r.
\end{aligned}
\]
Ordinary \ird{vy} supplies the transition from $\neg Q$ to $Q$; the
renewed-occurrence construction additionally requires the earlier $Q$.
Discourse reorientation by \ird{jas} alone supplies no such requirement.
For a state, these stages compare its validity over time. With a bounded
event, they compare contextually identified episodes or occasions, so a
contrary intervening stage does not deny that an earlier event occurred.
For departure, the intervening return resets the relevant situation; for
non-arrival, the successful intervening visit supplies the contrary occasion.
An arbitrary gap between two event times does not by itself identify this
discontinuity. The earlier occurrence and the contrary stage are background
conditions on this reading, rather than further events asserted by the
sentence; their projection under embedding and non-declarative force is not
independently specified.

Renewed occurrence also applies to a state or an ongoing process, where
\ird{vy} marks the attained validity of the predication without requiring
completion of the process. A gate can thus be open again after an intervening
closure, while progressive sleeping can resume after an interval of wakefulness.

\pex
\a \trsl{The western gate is open again.}
\lingcontext{English schema: the gate was open, closed, and has now reopened.}
\a
\begingl
\glpreamble \lingcontext{Janek was sleeping, woke for a while, and is now
sleeping once more.}//
\gla Janek jas vy zdrávime.//
\glb Janek.\Dir{} contrast already sleep-\Prog{}//
\glft \trsl{Janek is sleeping again.}//
\endgl
\a
\begingl
\glpreamble \lingcontext{Janek has been sleeping without interruption;
the speaker corrects an assumption that he has woken.}//
\gla Janek jas so zdrávime.//
\glb Janek.\Dir{} contrast still sleep-\Prog{}//
\glft \trsl{Janek is still sleeping, though.}//
\endgl
\xe

The state schema requires an earlier open gate, whereas the departure example
requires an earlier departure. A general renewed-occurrence reading therefore
does not identify a single sublexical attachment site. In the technical
repetitive--restitutive distinction, recurrence of a result state can be
licensed without an earlier event of bringing that state about: Nie's
Tagalog \textit{ulit} diagnostics distinguish a baby's earlier sleeping state
from an earlier event of falling asleep
\cite[§4.1, pp.~118--120, examples~47--50]{nie2020causatives}.
The established Iridian departure context supplies both an earlier event and
a reversed result, and consequently does not isolate a purely restitutive
interpretation of the verb. Lexical \ird{ru-} independently relates the event
to an earlier path or configuration: \ird{rustoj-} \trsl{retrace an earlier
route} specifies that relation even without \ird{jas vy}. A renewed occurrence
of going need not retrace the earlier route, while a progressive of
\ird{rustoj-} may describe retracing without completed return
(§\,\ref{sec:preverb-aspect}).

Additive perfective \ird{so} instead extends an established series by another
occurrence (§\,\ref{sec:class1-particles}). The two descriptions can apply to
the same sequence of events: a further arrival may also renew a previously
interrupted situation. Their licensing conditions differ, since membership
in a series supplies addition without itself specifying the contrary stage
required by recurrent \ird{jas vy}. Comparable distinctions between
continuation and addition of another event token are discussed for Totoli by
Bardají, whose particles nevertheless remain distributionally distinct even
where both permit an \trsl{again} translation
\cite[§§3.1, 3.3]{bardaji2025phasal}. Neither comparison restricts Iridian
renewed occurrence to states. With a habitual predication, recurrence of a
habit and repetition of one of its constituent episodes would supply
different comparison domains; the distribution of the renewed-occurrence
construction between those domains remains unspecified. Annoyance, weary
recognition, resignation and counterexpectation are contextual effects of
presenting a renewed situation, rather than additional encoded conditions.

Tagalog provides a close comparison at the level of the individual particles:
phasal \ird{vy} corresponds to \textit{na}, while discourse-reorienting
\ird{jas} corresponds to \textit{naman}. Independent \textit{naman} has
contrastive and non-contrastive uses, which AnderBois analyses through closure
of the preceding immediate question under discussion
\cite[§4]{anderbois2016naman}. That discourse operation alone introduces no
earlier event. McFarland's corpus study distinguishes recurrent
\textit{na naman} from ordinary \textit{naman} by its ordering: the former
remains together before interrogative \textit{ba}, although \textit{ba}
generally precedes independent \textit{naman}
\cite[pp.~1, 4]{mcfarland2001enclitics}. This cohesion provides a distributional
diagnostic for the Tagalog construction; Iridian adjacency follows from its
ordinary placement system and does not establish lexical fusion. The
distinction between a transition in an eventuality and a change
in discourse relevance is also maintained in Mazzitelli's comparative
analysis of change-of-state markers, which does not derive the recurrent
Tagalog combination \cite[§§2.1, 5]{mazzitelli2025aspect}.

Tagalog recurrence is compatible with imperfective processes:
\textit{Sumasakit na naman ang tiyan ko} \trsl{My stomach is aching again}
occurs among Jamindang's pain-predicate examples
\cite[pp.~88--89]{jamindang2022sakit}. Hsieh and Erlewine also document a
future-oriented recurrent use, but reject substitution of \textit{na naman}
for \textit{ulit} in their imperative meaning \trsl{Wash these clothes again}
\cite[§6.3, examples~36--37]{hsieherlewine2023nthu}.
These findings establish recurrence beyond completed events and a restriction
in the illustrated directive; they supply neither a general derivation of
\textit{na naman} nor Iridian judgements in corresponding environments.
In particular, the independent imperative uses of \ird{jas} and \ird{vy}
do not establish an iterative imperative \ird{jas vy}.

The combination with \ird{so} is less lexicalised and remains more transparently
compositional. Here \ird{so} contributes continuance, non-completion, or
ongoing validity, while \ird{jas} reorients that continuing situation against
some expectation, worry, competing line of action, or shared background
assumption. As a result, \ird{jas so} often yields translations such as
\trsl{still, though}, \trsl{still, after all}, \trsl{still, as you know}, or
\trsl{not yet, anyway}. This cluster is especially natural in concessive,
mitigating, reassuring, or regretive environments.

\pex
\a
\begingl
\glpreamble \lingcontext{Someone worries that there is no time left. The
speaker reassures them.}//
\gla Čas jas so jeca.//
\glb time.\Dir{} contrast still exist-\Cont{}//
\glft \trsl{There is still time, though.}//
\endgl
\a
\begingl
\glpreamble \lingcontext{The speaker softens a negative update.}//
\gla Tak jas so staly.//
\glb this.\Dir{} contrast still \Neg{}//
\glft \trsl{This is not the case yet, though.}//
\endgl
\xe

The same compositional logic extends to rhetorical interrogative environments.
When \ird{jas} combines with the interrogative particle \ird{by}, the result is often
best understood not as a neutral request for information but as interrogative
form plus discourse stance. The interrogative force is contributed by
\ird{by}; \ird{jas} then reorients the proposition against what should already
be evident, expected, or inferable in the current discourse. As a result,
\ird{by jas} often sounds retortive, incredulous, reproachful, or
challenge-like: roughly \trsl{come now, P?}, \trsl{P, after all?}, or
\trsl{was it not obvious that P?}, depending on context. The special flavour is
therefore pragmatic and compositional, not the fixed lexical meaning of an
idiom.

This same analysis also clarifies the larger clusters. \ird{By jas vy} combines
interrogative framing with an already-attained or now-salient result, and so
is especially natural where the speaker throws that result back into the
discourse as the obvious explanation for something else; this configuration
readily yields exasperated or \trsl{what else did you expect?} overtones.
\ird{By jas so}, by contrast, combines rhetorical interrogation with an open or
still-valid interval, and so more often suggests \trsl{is it not still the
case, after all?}, \trsl{is there not still time?}, or \trsl{has this not yet
been settled, then?} The contrast mirrors the broader difference between
\ird{jas vy} and \ird{jas so}: the former foregrounds attainment and, in
the renewed-occurrence construction, requires earlier validity and an
intervening contrary stage; the latter foregrounds persistence,
non-completion, or continuing relevance under a discourse-reorienting frame.
The first-arrival question in §\,\ref{sec:particle-combinations} illustrates
ordinary reorientation of attainment without recurrence. The availability
and background projection of a specifically recurrent question are not
established by that example.

With imperatives and hortatives, \ird{jas} often signals that the urged course
of action should already be evident from the shared situation. Depending on the
context, this may sound gentle, coaxing, impatient, teasing, or mildly
reproachful.

\pex
\begingl
\glpreamble \lingcontext{In a library, one speaker has been making noise for
some time. The other finally intervenes.}//
\gla Jas okouzé!//
\glb contrast be:noisy-\Proh{}//
\glft \trsl{Please do not be so noisy.}//
\endgl
\xe

\section{Evidential and stance particles}\label{sec:class3-particles}

Stance particles qualify the clause from outer positions in the particle
field, following its ordinary arguments and modifiers. The inventory comprises
\ird{oče}, \ird{nadě}, \ird{hladě},
\ird{kady}, and \ird{mlada}. Each may constitute an independent prosodic
word; ordinary initial stress remains available even when its audible
prominence is reduced in connected speech.

\subsection{Oče} \ird{Oče} is the mirative particle. It marks the proposition
as surprising, suddenly discovered, contrary to expectation, or otherwise
arriving at a moment when the speaker is not mentally prepared for it. The
relevant surprise may come from new information, abrupt realisation,
counterexpectation, or simple recollection of something momentarily forgotten.

\pex
\a
\begingl
\glpreamble \lingcontext{The speaker is surprised to learn that Janek has
left; they had expected him still to be there.}//
\gla Janek oče vy váržek.//
\glb Janek.\Dir{} mirative already leave-\Pf{}//
\glft \trsl{Janek has already left, unexpectedly.}//
\endgl
\a
\begingl
\glpreamble \lingcontext{The speaker suddenly remembers that it is raining.}//
\gla Oče pozebževí.//
\glb mirative rain-\Cont{}//
\glft \trsl{Oh, it's raining.}//
\endgl
\xe

\subsection{Nadě} \ird{Nadě} is the consequential particle. It presents the
proposition as the upshot of facts already in view, whether those facts have
just been stated overtly or are recoverable from the shared situation. English
translations often use \trsl{so}, \trsl{therefore}, or \trsl{that's why}, but
the particle is best analysed as a stance marker rather than as a simple
coordinator: it signals that the speaker is packaging the clause as a
conclusion that follows from what is already given.

\pex
\begingl
\glpreamble \lingcontext{The office is dark and Janek's coat is gone. The
speaker draws the obvious conclusion.}//
\gla Janek nadě vy váržek.//
\glb Janek.\Dir{} therefore already leave-\Pf{}//
\glft \trsl{So Janek has already left.}//
\endgl
\xe

\subsection{Hladě} \ird{Hladě} marks an inference grounded in some available
evidence while leaving room for uncertainty. It is more committal than open
speculation: the speaker presents the proposition as the likeliest conclusion
from what is known, but still as an inference rather than as directly verified
fact. It occupies the inferential position within Class~3, while
\ird{kady} marks less committed speculation. The inferential particle
\ird{hladě} has the broad pronunciation \ipagloss{[ˈxla.ɟɛ]}, with
\ird{dě} representing /ɟɛ/ under the general softness convention
(§\,\ref{sec:e-hook}).

\pex
\a
\begingl
\glpreamble \lingcontext{The speaker sees Janek's car outside the house and
infers that he has arrived, though they have not yet seen him in person.}//
\gla Janek hladě vy ščének.//
\glb Janek.\Dir{} infer already arrive-\Pf{}//
\glft \trsl{Janek must have arrived already.}//
\endgl
\a
\begingl
\glpreamble \lingcontext{The speaker has already seen Janek in person, so the
statement is no longer merely inferential.}//
\gla \ljudge{\#}Janek hladě vy ščének.//
\glb Janek.\Dir{} infer already arrive-\Pf{}//
\glft \trsl{Janek must have arrived already.}//
\endgl
\xe

\subsection{Kady} \ird{Kady} is the speculative particle. It presents the
proposition as conjectural, exploratory, or tentatively entertained rather
than as a conclusion anchored in clear evidence. It is especially natural in
questions or other non-assertive environments where the speaker invites the
addressee to consider a possibility without strongly committing to it.
The broad pronunciation is \ipagloss{[ˈka.dɨ̞]}, with initial stress and two
vowel nuclei. Initial /k/ is separated from /d/ by /a/, so there is no
obstruent cluster to trigger voicing; \ird{y} represents the central vowel
/ɨ/, with the standard lowered short realisation [ɨ̞], not the glide /j/
(§§\,\ref{sec:stress}, \ref{sec:lenition}, and~\ref{sec:orthographic-i} of
Chapter~\ref{ch:phon}).

\pex
\trsl{Might the oldest candle still be burning?}
\lingcontext{English schema: a tentative assessment of the candle's
continued burning, rather than a report or an inference from stated evidence.}
\xe

\subsection{Mlada} \ird{Mlada} is a hypothetical or desiderative particle
which frames the
clause as wished for, hoped for, counterfactual, or entertained against the
actual course of events. With conditional forms it yields ordinary
counterfactual protases and concessive hypotheticals; in isolation it may
function as an optative \trsl{if only} or a wish directed toward an uncertain
future outcome. Its meaning is therefore not merely temporal or modal in the
narrow sense, but explicitly speaker-oriented: the particle presents the whole
proposition as belonging to a non-actual, imagined, desired, or hoped-for
alternative.

\pex
\begingl
\glpreamble \lingcontext{His coming was already a desired possibility and remains one; Janek has not arrived.}//
\gla Janek mlada so ščének.//
\glb Janek.\Dir{} desiderative still arrive-\Pf{}//
\glft \trsl{I hope Janek still comes. / If only Janek would still come.}//
\endgl
\xe

In this desiderative construction, the input to \ird{so} is the availability
of the desired arrival, $Q$, rather than actual arrival, $P$. Its background
requires that same possibility $Q$ at the relevant preceding stage, not an
earlier actual arrival. \ird{Mlada} places the continuing possibility within
the wish; it does not generally cancel or filter presuppositions in its scope.

\section{Particle combinations, ordering and scope}
\label{sec:particle-combinations}

\ird{Janek oh vy ščének} combines addition with attained arrival, whereas
\ird{Janek jas vy váržek} places renewed departure within a contrasting
discourse frame. Both have adjacent particle sequences, although their
members retain distinct syntactic positions and semantic contributions.
Adjacency therefore does not establish a single lexical or phonological unit.
The three semantic classes distinguish families of meanings; their members
occupy the structural zones of §\,\ref{sec:th-particle-framework}, to which
those classes do not correspond one for one.

The principal structural relations are summarised below. The Comm-field is
a region containing heterogeneous operators, rather than an established
single \synhead{CommP}; ordinary placement and marked topic/focus
dependencies remain distinct from its operators
(Table~\ref{tab:th-clause-spine}).

\begin{table}[ht]
\footnotesize
\caption{Structural domains relevant to particle combinations.}
\label{tab:particle-combination-domains}
\begin{tabularx}{\textwidth}{l L L}
\toprule
Domain, highest first & Members & Contribution \\
\midrule
\synhead{SpeechActP} & \ird{nadě}, \ird{mlada} &
Consequential or desired/hypothetical framing \\
\synhead{EvidP} (outer) & \ird{oče}, \ird{hladě}, \ird{kady}, \ird{či} &
Mirative, evidential or epistemic stance \\
\synhead{PolP} & \ird{by}, or a null exponent & Interrogative clause typing \\
Comm-field & \ird{oh}, \ird{jas}, \ird{mó}, \ird{kamo}, \ird{dlí} &
Addition above reorientation; affirmation, exclusivity or interim priority \\
\synhead{PhasP} & \ird{vy} or \ird{so} &
Attainment or persistence over a polarised predication or licensed alternatives \\
\bottomrule
\end{tabularx}
\end{table}

Ordinary arguments and modifiers precede all these zones. Within the
field, outer material precedes \ird{by}, which precedes \ird{oh};
\ird{oh} precedes \ird{jas}, and \ird{jas} precedes \ird{vy} or
\ird{so}. This partial order does not determine every comment operator's
position relative to \ird{jas}, nor does it establish the co-occurrence of
two members of an outer zone. The phasal operators compete for one position
and cannot occur together, whereas unspecified combinations remain distinct
from prohibited ones (Table~\ref{tab:particle-combination-status}). Ordinary
arguments cannot interrupt the field. Ellipses, where used in schemata,
permit only independently licensed particle positions.

\begin{table}[ht]
\footnotesize
\caption{Distributional status of selected particle combinations.}\label{tab:particle-combination-status}
\begin{tabularx}{\textwidth}{L L}
\toprule
Combination & Status and restriction \\
\midrule
\ird{vy so}, \ird{so vy} & Excluded: a single phasal position \\
\ird{oh vy}, \ird{oh so}; \ird{mó vy}; \ird{dlí so} &
Licensed when both contributions and their backgrounds are satisfied \\
\ird{oh jas} &
Established discourse reorientation with addition \\
\ird{jas vy}, \ird{jas so} &
Established; recurrent \ird{vy} requires earlier validity and an intervening contrary stage \\
\ird{by jas vy/so} &
Established interrogative framing with discourse reorientation \\
\ird{kamo vy}, \ird{kamo so} &
Two bounded ordinary phasal compositions; restricted scalar low progress with \ird{so} \\
\ird{oh mó}, \ird{dlí kamo} &
Same-field co-occurrence and relative order unspecified \\
\ird{hladě kady}, \ird{oče či} &
Same-projection co-occurrence and relative order unspecified \\
\ird{hladě by jas so} &
Compatible with the zone sequence; the full combination is not independently specified \\
\bottomrule
\end{tabularx}
\end{table}

Distributional licensing and interpretive felicity impose separate
conditions. A permitted particle sequence can lack its intended reading
because the comparison domain is absent, the phasal background is false,
or the discourse does not supply the relation required by a higher operator.
Such a failure differs from the structural exclusion of \ird{vy so}.
Conversely, compatibility with the zone template supplies no co-occurrence
rule for a combination whose distribution is unspecified. Each licensed
reading must satisfy the semantic conditions of its construction.

\subsection{Co-occurrence within a structural field.}
The additive, affirmative, exclusive and interim meanings leave several distinct relations
unspecified. With \ird{oh} and \ird{kamo}, neither
co-occurrence nor a relation between the additive and exclusive alternative
sets is established: addition concerning one participant and exclusivity
concerning an action would differ from two operators associated with the
same participant set. With \ird{oh} and \ird{mó}, the unresolved
relation concerns addition and ratification of the proposition; with
\ird{dlí} and \ird{kamo}, it concerns the interim course
and the alternatives restricted by exclusivity. The corresponding
\ird{oh}--\ird{dlí}, \ird{mó}--\ird{dlí}, and
\ird{mó}--\ird{kamo} pairings likewise lack specified
co-occurrence and ordering rules. These are pair families rather than
licensed strings, and no general reversal or repetition rule follows from
their membership in the Comm-field.

The upper periphery contains two further kinds of unresolved combination.
Within outer \synhead{EvidP}, \ird{oče}, \ird{hladě}, \ird{kady},
and \ird{či} express surprise, inference, speculation, and attributed
information, whose possible joint realisation has no general specification.
Within \synhead{SpeechActP}, the same limit applies to combinations of
\ird{nadě} and \ird{mlada}; their combinations with lower \ird{jas}
also lack a general licensing rule. Across the outer projections,
the hierarchy supplies a relative structural height without independently
licensing simultaneous realisation. In particular, neither two meanings
being compatible in a discourse nor two particles belonging to different
projections establishes an Iridian construction. The grammar specifies no
general maximum number of particles: the exclusion of \ird{vy so}
concerns competition for the phasal position, whereas the unspecified
larger combinations concern co-occurrence rules.

Linear order, semantic nesting, and conventionalisation consequently remain
separate properties. Two licensed surface orders need not express two
scope relations, while a single order can support the two bounded
exclusive--phasal compositions below. Likewise, a recurrent contextual
translation does not establish a lexical unit: \trsl{finally} with
\ird{oh vy} retains its additive condition, and renewed occurrence with
\ird{jas vy} retains its earlier-occurrence and discontinuity
requirements without forming a single lexical unit. The
established \ird{Janek oh jas zdrávime} in the \ird{jas} entry combines
discourse reorientation with addition; it supplies no general permission
to insert further operators between those contributions.

\ird{Oh} ordinarily adds the complete phasal proposition to a relevant
alternative. When other arrivals have supplied the salient sequence and
Janek's arrival has long been outstanding, attained arrival can consequently
be understood as \trsl{finally}; prolonged waiting alone does not supply
addition. With \ird{so} and a negative arrival predicate, the added
proposition is instead continuing non-arrival.

\pex
\a
\begingl
\glpreamble \lingcontext{Other guests have arrived; after a long wait,
Janek arrives as well.}//
\gla Janek oh vy ščének.//
\glb Janek.\Dir{} also already arrive-\Pf{}//
\glft \trsl{Janek has finally arrived too.}//
\endgl
\a
\begingl
\glpreamble \lingcontext{Matuš has not arrived; Janek's non-arrival
likewise continues through the reference stage.}//
\gla Janek oh so zaščének.//
\glb Janek.\Dir{} also still \Neg{}-arrive-\Pf{}//
\glft \trsl{Janek has not arrived yet either.}//
\endgl
\xe

\ird{Mó vy} confirms the attained proposition, as in \ird{Janek mó vy
ščének} \trsl{Janek really has arrived}, while \ird{dlí so} places a
continuing situation within an interim plan. Thus \ird{Dá dlí so zdrava}
\trsl{For now I am still sleeping} requires a continuing sleep episode as
well as an interim course of action. Neither combination assigns the phasal
particle a new lexical meaning. Higher \ird{jas} similarly reorients the
phasal proposition, although its renewed-occurrence construction with
\ird{vy} adds a conventional earlier-occurrence condition:
\ird{Janek jas vy váržek}
\trsl{Janek has left again} requires an earlier departure and an intervening
return, rather than an attained departure alone. The detailed contrast with
more transparent \ird{jas so}, and with lexical restitutive
\ird{ru-}, together with the three-stage recurrent profile, belongs to the \ird{jas} entry in
§\,\ref{sec:class2-particles}.

Higher stance particles qualify the phasal content in their respective
ways. The English schema \trsl{Might the oldest candle still be burning?}
conjectures continued burning, while \ird{Janek nadě vy ščének}
\trsl{So Janek has already arrived} presents attained arrival as the
consequence under discussion. The continuative question invokes prior burning as background without
asserting that it continues now; the consequential assertion requires a
relevant discourse relation to preceding circumstances.

These combinations preserve distinctions between the content at issue and
the conditions under which it is presented. In ordinary phasal assertions,
the reference-stage predication is asserted against a presupposed preceding
stage; additive \ird{oh} requires a distinct relevant alternative, whereas
\ird{mó} ratifies the proposition and \ird{jas} relates it to the
discourse. Neither an English paraphrase nor the structural height of a
particle alone determines whether another attitude targets the predication,
its phasal background, or the act of presenting it. The formal scope
relations and their limits are set out in
§\,\ref{sec:th-particle-framework}.

Under interrogative \ird{by}, \ird{jas} reorients the phasal
proposition within the question; this does not require scope over the act
of asking. The retortive or reassuring effect depends
on the discourse; neither rhetorical force nor renewed occurrence follows
from this three-particle sequence alone. Interrogation leaves the current
attainment or persistence at issue, even when the context favours a positive
answer; it does not convert the phasal background into an assertion of that
answer. This ordinary interrogative composition does not determine the
licensing of a specifically recurrent reading under \ird{by}.

\pex
\a
\begingl
\glpreamble \lingcontext{Janek's first arrival explains why his coat is
here; the addressee continues to wonder about the coat.}//
\gla Janek by jas vy ščének?//
\glb Janek.\Dir{} \textsc{q} contrast already arrive-\Pf{}//
\glft \trsl{Has Janek already arrived, after all?}//
\endgl
\a
\begingl
\glpreamble \lingcontext{Time remains for the task, although the
addressee fears that it is too late.}//
\gla Čas by jas so jeca?//
\glb time.\Dir{} \textsc{q} contrast still exist-\Cont{}//
\glft \trsl{Is there still time, though?}//
\endgl
\xe

\ird{Kamo} precedes the phasal particle under both licensed ordinary
interpretations. In $\mathrm{EXCL}_{A}(\mathrm{PHAS}(P))$, exclusivity
compares complete phasal profiles; in
$\mathrm{PHAS}(\mathrm{EXCL}_{A}(P))$, the exclusive proposition itself
is attained or maintained. The western-gate schemata in
§\,\ref{sec:th-exclusive-phasal} distinguish \trsl{Only the western gate
is still open}, which permits a newly opened gate, from \trsl{It is still
only the western gate that is open}, which requires exclusivity at both
comparison stages. With fixed alternatives and the same two stages,
continuing exclusivity entails exclusivity of persistence, so only the first
reading can be isolated by a context in which another gate has just opened.

Attainment permits separation in either direction. In the following English
schemata, the alternatives are the same three gates and the predicate is
being open:
\pex
\a \trsl{Only the western gate is already open.}
\lingcontext{The eastern gate was open and remains open; the western gate
has just opened and the northern gate is closed. Only the western gate
satisfies the complete transition-to-open profile.}
\a \trsl{It is already only the western gate that is open.}
\lingcontext{The western and eastern gates were open; the eastern gate
has closed and the northern gate remains closed. Exclusivity is newly
attained, although the western gate was open before.}
\xe
The first context isolates exclusivity over attainment, while the second
isolates attainment of exclusivity. An alternative lacking the preceding
state fails the relevant complete phasal profile without making the whole
exclusive comparison undefined. These constructions operate within one finite
domain without intervening Mood or clausal evidentiality; they supply no
general lowering rule for other comment-domain particles. Their formal
definitions and polarity diagnostics are given in
§§\,\ref{sec:th-exclusive-phasal}--\ref{sec:th-polarity-diagnostics}.

Scalar \ird{kamo} under temporal \ird{so} additionally permits a
restricted low-progress interpretation. The English schemata
\trsl{The scribe has copied only three folios so far} and
\trsl{The climbers have reached only the lower ridge so far} compare an
attained stage with an earlier stage and a salient higher goal. Unlike
persistence of an entire exclusive proposition, this construction preserves
the upper bound at the earlier stage while anchoring MIN at the reference
stage. The scribe may therefore have copied one folio earlier and three now,
without having reached three earlier. The ordered progression must contain
a genuine advance and a higher relevant attainable stage; further advance
remains a defeasible expectation. Temporal focus particles can share a
semantic core while selecting distinct
alternative structures, as in the analysis of Dutch by Neeleman and van de Koot
\cite[§§2--3]{neelemanvandekoot2022}. The construction and its separate
background conditions are specified in §\,\ref{sec:th-exclusive-phasal}.

Related English schemata vary the measured progression and the diagnostic
scope relation:
\begin{table}[ht]
\footnotesize
\caption{Supplementary phasal and progress schemata.}
\label{tab:kamo-phasal-schemata}
\begin{tabularx}{\textwidth}{L L}
\toprule
English schema & Interpretation or diagnostic \\
\midrule
Only the oldest candle is still burning. & Exclusivity of persistence; a newly lit candle is permitted. \\
It is still only the oldest candle that is burning. & Persistence of exclusivity over the same candles. \\
Only the eastern road is not flooded yet. & Exclusivity over persistence of a negative state. \\
The floodwater has risen only to the second step so far. & A low degree can be reassuring; disappointment is not entailed. \\
The carpenter has repaired only the doorframe so far. & Incremental completion within a larger project. \\
The child has learned only the first stanza so far. & Sequential attainment of text portions. \\
Only one bell has stopped ringing already. & A diagnostic for exclusivity with attainment of a change of state. \\
\bottomrule
\end{tabularx}
\end{table}
The bell schema does not independently license a new scope reading, just as
the progression schemata do not determine their eventual aspect morphology.

Temporal \ird{so} in this construction differs from scalar-additive
\ird{so} \trsl{even}. In the English diagnostic
\trsl{Even the master mason repaired merely the doorway}, addition would
compare actors, whereas the upper bound would compare repair accomplishments:
$C_{\mathrm{so}}$ and $C_{\mathrm{kamo}}$ need not coincide. If both
operators instead compare the same alternatives, an additive requirement
for another true alternative can conflict with the exclusive upper bound.
Neither semantic compatibility nor the shared pronunciation \ird{kamo so}
independently licenses this scalar-additive combination. Likewise,
\trsl{Only the harbourmaster received another warning} tests an additive
event interpretation of \ird{so}, for which the temporal scope rules alone
supply no constructional licence.

Tagalog supplies a comparison in which combination, linear order and
available scope are independently constrained. Table~\ref{tab:tagalog-combination-scope}
summarises Hsieh and Erlewine's tested contrasts within one clitic domain
\cite[§§3--4.2, appendix~A]{hsieherlewine2024}. The reversed \textit{lang na}
and \textit{lang pa} orders are excluded there; by contrast, \textit{na lang}
and \textit{pa lang} are permitted, although the reverse scope reading is
infelicitous in their diagnostic contexts. Variation in the order of
\textit{din} and \textit{lang} likewise does not yield scope permutation.
These findings establish no general freedom to stack particles.

\begin{table}[ht]
\footnotesize
\caption{Selected Tagalog clitic combinations: order and interpretation.}
\label{tab:tagalog-combination-scope}
\begin{tabularx}{\textwidth}{L L}
\toprule
Surface combination & Scope supported by the cited diagnostics \\
\midrule
\textit{na lang}, \textit{pa lang} & PHAS $>$ EXCL; EXCL $>$ PHAS unavailable \\
\textit{lang din}, \textit{din lang} & ADD $>$ EXCL in either order, with ordering preferences \\
\textit{na rin} & Co-occurrence established; relative ADD/PHAS scope not settled by the tested readings \\
\textit{hindi na/pa/rin} & PHAS/ADD $>$ NEG despite negation's precedence \\
\textit{hindi lang} & Both NEG $>$ EXCL and EXCL $>$ NEG \\
\bottomrule
\end{tabularx}
\end{table}

The restriction to one clitic domain matters: the same study's topic and
focus constructions with separate domains permit the converse relation
between addition and exclusivity \cite[§4.2, fn.~5]{hsieherlewine2024}.
Their 2025 draft handbook retains this distinction between clitic domains
\cite[§4.2]{hsieherlewine2025handbook}. Nor does co-occurring \textit{na rin} establish addition over the whole
phasal proposition: presupposition projection limits that diagnostic
\cite[appendix~A]{hsieherlewine2024}. Iridian's full-profile additive
requirement and its two exclusive--phasal readings are therefore independent
constructional rules, rather than consequences of the Tagalog comparison.

Hsieh and Erlewine's analysis of low-progress \textit{pa lang}
uses an ordered progression: the reached stage is neither initial nor final,
and continued progress is a cancellable expectation
\cite[§4, pp.~5--7]{hsieherlewine2023nthu}. The handout's derivation maintains an upper
bound over the progression at the earlier stage, while the lower-bound
presupposition contributed by \textit{lang} is evaluated at the current
stage. It consequently need not presuppose that the presently reached stage
held earlier. Iridian's restricted low-progress construction preserves this
separation, whereas ordinary persistence applies to the complete exclusive
proposition. The change-of-plan \textit{na lang} account relates a less desirable
replacement plan to a covert metaphysical necessity modal and a director
capable of, and committed to, ensuring its realisation. The handout assumes
exclusivity above the modal while explicitly leaving their relative scope
uncertain \cite[§5, pp.~7--9]{hsieherlewine2023nthu}.
These analyses illustrate how construction-specific licensing can emerge
from composition, without deriving Iridian renewed occurrence or supplying
its missing co-occurrence rules.

Upper-particle combinations introduce a further distinction between the
source of a report and the holder of an inference or speculation. Tan
Almazán describes Tagalog \textit{daw yata} as permitting a report of an
inference, while the more restricted reverse order can express an inference
based on a report; acceptance of that reverse order varies among the
consultants \cite[§3.3.2.2.2, pp.~127--128]{tanalmazan2019}. By contrast,
fixed \textit{daw kayâ} permits both reporting another person's wondering
and wondering about reported content
\cite[§3.3.2.2.3, pp.~128--129]{tanalmazan2019}. Neither the number of
surface orders nor a general reportative label therefore determines the
available nesting. These contrasts clarify the distinct source and attitude
relations left unspecified for Iridian's outer-particle pairings; they
do not establish corresponding Iridian readings.

Other comparisons reinforce the need to diagnose each reading separately.
Sentence-final Singlish \textit{already} has clause-wide scope despite the
clause-medial position of related Chinese particles
\cite{erlewine2023already}. Beck places reaffirmative German \textit{noch}
and English \textit{still} above an assertion operator, whereas her
order-of-mention reading concerns German \textit{noch} above an answer
operator \cite[§4.3, pp.~157--159]{beck2016still}. Neither comparison places
ordinary continuative Iridian \ird{so} in the speech-act domain.

The distinction between a field and a fully specified sequence also has
comparative support outside Tagalog. Lohmann and Koops's study of spoken
German separates ordering preferences from the strength of co-occurrence
associations; some particle pairs occur systematically in both orders, so
a single ordering hierarchy does not capture their combinatorial behaviour
\cite[§§4--5.3]{lohmannkoops2026}. Frequent combination does not by itself
establish a non-compositional lexical unit
\cite[§5.4]{lohmannkoops2026}. These findings support keeping order,
co-occurrence, and conventionalisation distinct, without supplying either
reversibility or frequency-based judgements for Iridian.

Fang and Hengeveld's analysis of Mandarin sentence-final particles further
distinguishes membership in a semantic category from membership in a scope
layer: particles in different categories can combine even when assigned to
the same layer \cite[§5.5]{fanghengeveld2022}. Their Functional Discourse
Grammar layers are not Iridian's Minimalist projections, but the distinction
clarifies why a shared \synhead{EvidP} label cannot decide the status of
\ird{oče} with \ird{či}, just as a common stance function cannot by
itself establish redundancy. Iridian's same-projection combinations remain
unspecified independently of this comparison.

Negation supplies the phasal input whether it is realised by \ird{za-},
by suppletive \ird{zad-}, or by equative \ird{staly}. The ordinary readings
are consequently \ird{vy}$(\neg P)$ \trsl{no longer} and
\ird{so}$(\neg P)$ \trsl{not yet}, rather than
$\neg\mathrm{VY}(P)$ or $\neg\mathrm{SO}(P)$. In particular,
\ird{Marek doktor so staly} and \ird{Pivo so zada} retain this scope
although the pronounced negatives differ in position and morphological
status (§\,\ref{sec:class1-particles}). Asp.Align supplies the viewpoint
and alignment of a verbal predicate before phasal interpretation. An
additive reading with a completed-event predicate requires an independently
established series, while progressive predicates can supply an ongoing
state. Contemplative and desiderative
combinations concern a projected or available possibility rather than an
asserted completed event. Interrogation takes the resulting phasal
proposition within its scope, as in the \ird{jas} questions above.

Directive and irrealis combinations preserve the same separation of phasal
meaning from clause type. \ird{Mó vy pjaštím!} \trsl{Go on, eat already!}
urges initiation without asserting completed eating. Negative Mood supplies
its licensed negative input below Phas, so \ird{Oh so pjašté} calls for
continued restraint as an additional course of action. The morphology and
interpretation of these negative non-indicatives are specified in
§\,\ref{sec:vm-moodp}, with their controlled contexts in
§\,\ref{sec:th-polarity-diagnostics}; the pronounced Mood suffix does not
place negation outside the phasal operator.

Prosodic dependence follows the local host rather than the full semantic
scope of the combination. In \ird{Janek jas vy váržek}, \ird{vy} leans
on \ird{váržek} without shifting its lexical stress or contributing a
syllable to its metrical parse, while \ird{jas} retains a separate position
and need not share the predicate's stress domain.
Comment-domain particles may be weakly attached to adjacent following
material, and outer stance particles permit independent initial stress.
Phrasal deaccentuation can reduce prominence without changing either this
lexical stress or the syntactic position (§\,\ref{sec:stress}). No general
rule assigns one stress domain to all co-occurring particles, nor does the
structural hierarchy alone specify a pause or an intonational boundary at
every zone edge. Attribution and embedding can introduce distinct domains
for a particle's contribution
(§\,\ref{sec:particle-domains-attribution}); clause-final discourse particles
have a separate utterance-edge distribution whose interaction with the
preverbal series is described in
§\,\ref{sec:final-particle-combinations}.

\section{Structural analysis of the particle field}
\label{sec:th-particle-framework}

The adverbial particle system, analysed in detail in §\,\ref{sec:adv-particles}
establishes the main empirical evidence for
the left-peripheral structure summarised in Figures~\ref{fig:th-clause-fields}
and \ref{fig:th-left-periphery}. The particles serve here as syntactic
diagnostics: their ordering across structural zones, their scope properties,
and their ability or inability to license fragment answers jointly constrain
the structural analysis. Surface position alone does not
determine every scope relation, particularly the restricted lower exclusive
interpretation described in §\,\ref{sec:th-exclusive-phasal}.

Three adverbial-particle zones and a separate clause-typing position are
 distinguished within the finite clause. Ordinary arguments and modifiers
 precede all three zones and the clause-typing head through the dependency
 specified in §\,\ref{sec:pivot-placement}; independently linked topics and
 marked focus precede that ordinary placement region.

\begin{description}[leftmargin=2em, labelwidth=!, labelindent=0em]
\item[\upshape Clause typing (\synhead{PolP}).] The interrogative particle
 \ird{by} realises locally head-initial Pol\textsuperscript{0}, below outer
 evidential and speech-act projections but above the comment operators.
 Its complement follows it, while arguments pronounced at the finite
 placement edge precede it (§\,\ref{sec:by-pol-head}). Pol names the
 interrogative operator of this analysis; its position below marked focus and
 ordinary placement is specific to the Iridian structure. Cartographic Force,
 Fin and Int distinguish other clause-edge functions and positions, including
 an Italian Int above Foc \cite[§§2.2--3.1]{rizzibocci2017}. Those labels do
 not independently establish the position or interpretation of Iridian \ird{by}.
\item[\upshape Zone~\textsc{i} (outer evidential and speech-act).]
 Speech-act particles such as \ird{nadě} and \ird{mlada}, and evidential
 and epistemic particles such as \ird{oče}, \ird{hladě}, \ird{kady} and
 \ird{či}, occupy articulated positions above \synhead{PolP}. Their
 semantic relations to propositions, sources and questions remain distinct.
 Their pronounced positions follow ordinary constituents, including those
 without a topic interpretation.
\item[\upshape Zone~\textsc{ii} (comment-domain).] The Comm-field contains
 additive \ird{oh}, affirmative \ird{mó}, exclusive \ird{kamo},
 interim-sequential \ird{dlí}, and reorienting \ird{jas}. Within the
 specified partial order, \ird{oh} precedes \ird{jas}, which precedes
 phasal particles; \ird{by} precedes both. Comm names an articulated region,
 rather than a single head or a uniform semantic type. No general ordering
 or co-occurrence rule follows for its other members.
\item[\upshape Zone~\textsc{iii} (phasal).] The phasal particles \ird{vy}
 (\trsl{already}) and \ird{so} (\trsl{still} / \trsl{not yet}) occupy
 \synhead{PhasP} above the polarity-bearing predication: indicative
 \synhead{NegP}, when present, otherwise the viewpoint/lexical predicate
 domain. Their input is a positive or negative proposition or a contextually
 ordered alternative licensed in §\,\ref{sec:class1-particles}; they do not
 replace Asp.Align.
\end{description}

\noindent The ordinary surface order is arguments and modifiers $>$ outer
 evidential/speech-act material $>$ \ird{by} $>$ \ird{oh} $>$ \ird{jas}
 $>$ \ird{vy}/\ird{so} $>$ verbal predicate. This sequence states the
 relations among named positions, without licensing every combination or
 settling the order of every stance and comment operator. Arguments do not
 interrupt the ordinary particle field. Clause-final particles retain their
 utterance-edge positions, while nonverbal predication has the independently
 specified structures of Chapter~\ref{ch:predicates}.

The ordering is therefore a partial ordering over independently licensed
positions, rather than a template into which every member of a semantic class
can be inserted. Within a shared field, complementary occupation of one
position, distinct ordered subpositions, and restricted adjunction would make
different predictions about co-occurrence and order; the unspecified Iridian
pairs do not yet distinguish these analyses. The exclusion of \ird{vy so}
and \ird{so vy} follows from their established competition for one phasal
position, whereas no corresponding single-occupant rule has been established
for the Comm-field or outer \synhead{EvidP}. Conversely, the compatibility of
two independently described meanings does not establish their joint syntactic
licensing. Pairwise licensing alone does not establish a sequence of three
particles.

Interpretation further depends on the object to which each contribution
applies. A phasal operator relates a polarised predication or another licensed
input to preceding and reference stages; a focus-associated operator compares
alternatives whose domain and associate must be identified; an evidential or
stance particle relates its content to a source, a perspective holder, or a
discourse relation. Identifying the source of a report separately from the
holder of an assessment prevents a single perspective from being assigned to
several operators merely because they share a field. Tagalog evidential
combinations illustrate the relevance of this distinction
\cite[§3.3.2.2]{tanalmazan2019}. Likewise, a proposition, a
question concerning that proposition, and the act of asking that question
are distinct possible objects of qualification. The upper-peripheral position
of an Iridian particle establishes its general relation to the lower clause
without, by itself, choosing among these more specific objects or supplying
an otherwise unspecified source or holder. The distinction between a local
syntactic domain and the attribution of its content is developed for embedded
clauses in §\,\ref{sec:particle-domains-attribution}.

Semantic scope must also be distinguished from the projection of background
content. When a phasal or additive background remains available under
interrogation, that persistence does not itself show that its introducing
particle takes scope over the question. The at-issue contribution and the
background conditions must be traced separately through the composition,
with the relevant stages and alternatives held constant. Hsieh and Erlewine's
Tagalog diagnostics illustrate both the value of this separation and its
limits: the projected presuppositions of \textit{na rin} can leave competing
scope analyses indistinguishable \cite[appendix~A]{hsieherlewine2024}.
Iridian's specified exclusive--phasal rules below consequently distinguish
interpretations by their complete profiles, rather than deriving scope from
linear precedence or from background persistence alone.

The renewed-occurrence construction with reorienting \ird{jas} and
phasal \ird{vy} likewise has a specified background: earlier validity of
the polarised predication precedes the contrary stage against which its
current attainment is asserted (§\,\ref{sec:class2-particles}). That
three-stage profile supplements ordinary \ird{vy}; neither the structural
height of \ird{jas} nor its discourse-reorienting contribution alone
introduces the earlier occurrence. The construction retains the
distinct positions of \ird{jas} above \ird{vy}, with ordinary constituents
preceding both. Its declarative profile does
not specify projection of the earlier-occurrence condition through Mood,
interrogation, or clausal embedding.


\section{Phasal scope and focus-associated exclusivity}
\label{sec:th-exclusive-phasal}

The lower interpretive order is comment-domain operators $>$ phasal operator
$>$ polarised predication $>$ viewpoint aspect and the inner lexical predicate. Thus the
ordinary negative combinations are $\mathrm{VY}(\neg P)$ and
$\mathrm{SO}(\neg P)$, rather than negation of an already/still proposition.
The distinction is independent of whether negation is pronounced as a prefix,
as existential suppletion, or as clause-final \ird{staly}. The relation between
phasal interpretation and the realisation of finite heads is compositional,
not a requirement to reorder the verbal suffixes.

In negative non-indicatives, a silent interpretable $\mathrm{Op}_{\neg}$
applies to the AspP proposition below Phas, while the higher Mood head bears
its concord feature and selects the negative suffix. The clause-local
licensing and composition are specified in §\,\ref{sec:vm-moodp}:
$M(C(\mathrm{PHAS}(\neg P)))$. Thus a prohibitive with \ird{so} calls for
continued restraint; it does not compose as $\neg\mathrm{SO}(P)$ and does
not introduce indicative NegP or \ird{za-}. The suffix's pronounced position
is not its semantic scope. The same lower polarity input is used under
hortative and irrealis Mood.

Focus-associated \ird{kamo} has two specifically licensed
compositions with \ird{vy}/\ird{so}. Its ordinary comment-domain interpretation
is $\mathrm{EXCL}_{A}(\mathrm{PHAS}(P))$, where $A$ is the contextually fixed
set of focus alternatives. A lower interpretation
$\mathrm{PHAS}(\mathrm{EXCL}_{A}(P))$ is available when that entire exclusive
proposition supplies the attained or continuing state. The focus associate
and alternative set must remain constant across the relevant stages.
Precisely, let $p_x(t)$ be the proposition with the focused associate replaced
by $x$, with actual associate $a\in A$ and all other material fixed.
For argument focus, $x$ ranges over individuals; for VP focus it ranges over
contextually comparable actions, with their participants held fixed. For
truth-valued alternatives $q_x$, write
$E_A(q)=q_a\wedge\bigwedge_{b\in A\setminus\{a\}}\neg q_b$.
The bounded lower rule constructs the stagewise proposition
$q(t)=E_A(x\mapsto p_x(t))$ and supplies that proposition to PHAS.
If the prejacent is negative, $p_x(t)$ already contains the one licensed
negation; the rule neither duplicates it nor moves it over PHAS.

In the ordinary high reading, exclusivity instead compares complete phasal
profiles. With the same reference stage $r$ and preceding stage $r^{-}$ for
all alternatives, define the truth-valued profile predicates
\[
\begin{aligned}
h_{\mathrm{VY}}(x)&=\neg p_x(r^{-})\wedge p_x(r),\\
h_{\mathrm{SO}}(x)&=p_x(r^{-})\wedge p_x(r).
\end{aligned}
\]
The high construction requires $h_{\mathrm{PHAS}}(a)$ and excludes
$h_{\mathrm{PHAS}}(b)$ for every distinct $b\in A$. The actual associate
retains its presupposed preceding-stage background, while each excluded
alternative is tested for satisfaction of the complete two-stage profile.
An alternative that lacks the relevant preceding state therefore fails that
profile without making the high exclusive construction undefined. Thus a
newly opened gate does not satisfy the profile of continued openness,
while a gate that was already open does not satisfy the attainment profile. The scope notation
$\mathrm{EXCL}_{A}(\mathrm{PHAS}(P))$ abbreviates this construction-specific
comparison; it does not negate a partial phasal denotation whose
presupposition has failed. These profile predicates formalise the ordinary
two-stage approximation, without replacing the phasal presupposition of the
actual associate or extending the rule to scalar constructions. For a negative
prejacent, $p_x$ remains the already polarised proposition in both profiles.

Only the \ird{kamo}--\ird{vy}/\ird{so} construction licenses the
lower exclusive composition, within one finite domain and without intervening Mood or
clausal evidentiality. It does not license lowering \ird{oh}, \ird{mó},
\ird{dlí}, or an arbitrary sequence of comment operators.

The particle remains pronounced in the Comm-field above the phasal particle.
A topicalised focus associate binds its lower argument
variable $x$ (equivalently, its interpreted lower copy supplies that slot);
the pronounced topic itself is not moved into the phasal complement at LF.
Focus association requires a contrastive set containing that referent, even
when the same nominal is the aboutness topic. The lower rule reconstructs the
exclusive proposition, not all high material. These two rules are specified
Iridian conventions. Erlewine's lower-copy mechanism concerns English
\textit{even}; his discussion explicitly does not extend that mechanism to
English \textit{only} \cite[p.~863, fn.~9]{erlewine2016copies}. It therefore
neither derives Iridian exclusivity nor establishes backward association to a
topic. Here topic association is explicitly variable binding, and the lower
alternative-proposition rule is stipulated for the Iridian construction.

The following English schemata specify interpretations and discourse
conditions; they supply neither Iridian surface forms nor morphological
analyses.

\pex
\a \trsl{Only the western gate is still open.}
\lingcontext{$\mathrm{EXCL}_{A}(\mathrm{SO}(P))$. The western and eastern
 gates were open yesterday; the eastern gate has closed, while a third gate
 has just opened. Only the western gate satisfies the continuing profile.}
\a \trsl{It is still only the western gate that is open.}
\lingcontext{$\mathrm{SO}(\mathrm{EXCL}_{A}(P))$. The western gate was the
 sole open gate yesterday and remains the sole open gate today.}
\xe

The lower reading requires prior exclusivity, whereas the high reading
requires the western gate to be the sole continuing member of the earlier
set of open gates; a newly opened alternative need not have been open before.
With the same gate alternatives at both stages,
$\mathrm{SO}(\mathrm{EXCL}_{A}(P))$ entails
$\mathrm{EXCL}_{A}(\mathrm{SO}(P))$, although the converse fails. The
first context therefore excludes the lower reading, while the second supports
it without independently ruling out the high reading. The contrast concerns
the stronger requirement of continuing exclusivity, rather than two
mutually exclusive contextual distributions. The corresponding candle
schemata, \trsl{Only the oldest candle is still burning} and
\trsl{It is still only the oldest candle that is burning}, admit the same
comparison: a newly lit candle is compatible with exclusivity of continued
burning, whereas another current flame contradicts continuing exclusivity.

\noindent\textit{Scalar alternatives and low progress.}\enspace
\ird{Kamo} also bounds alternatives ordered by degree or contextual rank,
with the MIN/MAX interpretation defined in §\,\ref{sec:class2-particles}.
The Boolean $E_A$ above is a satisfaction diagnostic for the ordinary
individual- or action-exclusive cases, in which the prejacent is true and
all distinct relevant alternatives are false; it does not define every
scalar use or identify the asserted and presupposed contributions. In
\trsl{The child can read only the simplest passages}, the relevant alternatives
concern attainable difficulty levels under the ability modal. Several easy
passages may be readable, since the upper bound excludes ability beyond the
specified level rather than every other passage. This is a schematic
interpretation of scalar \ird{kamo}; the syntax of its association within
an overt modal construction requires the independently licensed modal
structure. Contextual ordering, rather than an unrestricted choice of
lexical senses, distinguishes this use from individual exclusivity
\cite[§2.3]{coppockbeaver2014}.

Low-progress \ird{kamo so} has a narrower interpretation than persistence
of an entire exclusive proposition. It describes an attained stage below a
salient expected goal while an upper bound continues from an earlier stage.
The lower bound is evaluated at the reference stage, so an earlier stage may
fall below it. This construction extends the ordinary lower interpretation
within the same finite phasal domain; it does not make every scalar use of
\ird{so} interchangeable with temporal persistence. The separation of
current attainment from a continuing upper bound follows the comparative
analysis of Tagalog \textit{pa lang}, while its licensing in Iridian is a
constructional convention \cite[§4]{hsieherlewine2023nthu}.

Let $d(t)$ measure cumulative attainment on a contextually ordered scale,
$k$ be the designated attained threshold, $k_{\mathrm{next}}$ the next
relevant higher threshold, and $g_C$ a salient expected goal. For
$r^{-}<r$, the specialised declarative profile is:
\[
\begin{aligned}
\text{current lower bound (MIN):}\quad & d(r)\geq k,\\
\text{asserted upper bound (MAX):}\quad & d(r)<k_{\mathrm{next}},\\
\text{preceding upper bound:}\quad & d(r^{-})<k_{\mathrm{next}},\\
\text{non-vacuous progression:}\quad & d(r^{-})<d(r).
\end{aligned}
\]
The current lower bound and preceding upper bound are background conditions;
the progression condition restricts this construction to an actual advance
between the chosen stages. The ordering must make $k<g_C$ contextually
salient, without entailing that the goal will be reached or that the speaker
regrets the limited progress. On a dense scale without a next threshold,
MAX instead requires that every true relevant attainment alternative be
ranked no higher than $k$. These conditions leave ordinary \ird{so} and
ordinary persistence of exclusivity unchanged.

\pex
\a \trsl{The scribe has copied only three folios so far.}
\lingcontext{One folio was complete at the earlier stage; three are complete
 now, against an expected twenty. The continuing upper bound is fewer than
 four completed folios; attainment of three is required only now.}
\a \trsl{The climbers have reached only the lower ridge so far.}
\lingcontext{The party has advanced from the trailhead to the lower ridge;
 the upper ridge and summit remain higher attainable milestones.}
\xe

The scribe's earlier stage satisfies MAX but not the current MIN of three;
requiring the entire MIN/MAX package to hold earlier would therefore reject
this progression. If three folios had already been complete and no further
progress had occurred, ordinary persistence of exclusivity could hold,
although the non-vacuous low-progress profile would fail. Cardinal quantities
and ordered geographical milestones consequently instantiate the same
construction without requiring the scale to count objects. Temporal focus
particles can impose distinct requirements on their alternative structures,
as the Dutch comparison shows \cite[§§2--3]{neelemanvandekoot2022}. In the additional
schema \trsl{The floodwater has risen only to the second step so far}, limited
progress may be reassuring; disappointment is no part of the upper bound.

Temporal \ird{so} in these profiles differs from scalar-additive \ird{so}
\trsl{even}, whose alternatives require addition and a contextual ordering
(§\,\ref{sec:class1-particles}). An exclusive upper bound and an additive
witness cannot be imposed indiscriminately on the same alternatives. In
\trsl{Even the master mason repaired merely the doorway}, the intended
scales concern the actor and the extent of the repairs respectively:
$C_{\mathrm{so}}$ and $C_{\mathrm{kamo}}$ must be specified separately.
This is a diagnostic schema, not an additional licensed Iridian particle
combination. Likewise, \trsl{Only the harbourmaster received another warning}
tests exclusivity over event addition, while \trsl{Only one bell has stopped
ringing already} tests a change-of-state prejacent under \ird{vy}; neither
English schema settles an unspecified Iridian co-occurrence or association
rule.

MIN, MAX and phasal backgrounds must be tracked separately under \ird{by},
negative Mood and epistemic qualification. The declarative low-progress
profile does not determine whether its progression condition projects under
those operators. Questioning MAX while retaining MIN and negating MAX while
retaining MIN are distinct formal diagnostics; their availability does not
establish those projection patterns as Iridian rules. The following bounded
paradigm retains the established polarity and phasal conventions, while the
extended semantic diagnostics distinguish lower-bound failure, upper-bound
failure and absence of progression.


\section{A bounded polarity--phase--focus diagnostic paradigm}
\label{sec:th-polarity-diagnostics}

The following constructed paradigms specify Iridian design judgements, not
attested speaker responses. They hold the ordinary two-stage interpretation
fixed. Write $V=\mathrm{VY}$, $S=\mathrm{SO}$,
$A=\mathrm{ADD}_A$ for additive \ird{oh}, and $E=\mathrm{EXCL}_A$ for
\ird{kamo}; $A$ in the subscript denotes a fixed alternative set.
$A(q)$ asserts $q_a$ and presupposes some distinct $q_b$, whereas $E(q)$
checks that $q_a$ is true and every distinct $q_b$ in that set is false.
This satisfaction check abstracts away from the MIN/MAX division of content. In the tables,
$E(V(P))$ and $E(S(P))$, including negative prejacents and embedding under
Mood, abbreviate the high profile comparison defined in
§\,\ref{sec:th-exclusive-phasal}; $E(P)$ inside $V$ or $S$ denotes
stagewise truth-valued exclusivity. Additivity here associates with the whole phasal proposition, not an independently lowered
operator. Its additional alternative must satisfy both the preceding-stage
background and the reference-stage condition of that phasal profile; an
alternative whose background fails cannot supply the additive presupposition.
The associate is \ird{Marek} in Table~\ref{tab:th-indicative-diagnostics}:
$p_x(t)$ means that $x$ is sleeping at stage $t$. The topic binds its lower
argument variable as in §\,\ref{sec:th-exclusive-phasal}.

\begin{table}[ht]
\footnotesize
\caption{Indicative diagnostics: fixed person alternatives.}
\label{tab:th-indicative-diagnostics}
\begin{tabularx}{\textwidth}{L l l}
\toprule
\textsc{surface string} & \textsc{ordinary scope} & \textsc{status} \\
\midrule
\ird{Marek oh vy zdrávime.} & $A(V(P))$ & L \\
\ird{Marek oh so zdrávime.} & $A(S(P))$ & L \\
\ird{Marek kamo vy zdrávime.} & $E(V(P))$ & L \\
\ird{Marek kamo so zdrávime.} & $E(S(P))$ & L \\
\ird{Marek oh vy zazdrávime.} & $A(V(\neg P))$ & L \\
\ird{Marek oh so zazdrávime.} & $A(S(\neg P))$ & L \\
\ird{Marek kamo vy zazdrávime.} & $E(V(\neg P))$ & L \\
\ird{Marek kamo so zazdrávime.} & $E(S(\neg P))$ & L \\
\bottomrule
\end{tabularx}
\end{table}

Here L means licensed with its stated background, not felicitous in every
context. For each asserted alternative, $V(P)$ requires non-sleeping before
and sleeping now; $S(P)$ requires sleeping at both stages; $V(\neg P)$
requires sleeping before and non-sleeping now; $S(\neg P)$ requires
non-sleeping at both stages. In an \ird{oh} row at least one other person
must satisfy the same complete phasal description; in a high \ird{kamo} row
none may do so. A background mismatch for the actual associate, or the
absence of a suitable additive witness in an \ird{oh} row, makes the
intended ordinary reading unavailable; an excluded exclusive alternative
need not satisfy that background. For the \ird{kamo} rows only, the alternative lower reading is
$V(E(P))$, $S(E(P))$, $V(E(\neg P))$, or $S(E(\neg P))$, respectively,
with the entire exclusive proposition evaluated at both stages. This is
licensed only by the bounded rule above. In particular, only Marek being
\emph{still not asleep} need not imply that everyone else is asleep now;
\emph{still only Marek not asleep} does imply that, and also prior exclusivity.

For the matched negative non-indicatives below, $P$ is eating with agent and
food fixed by context. Focus is the VP \trsl{eat}, with a fixed set of
alternative actions such as eating and drinking; it is not a free change of
subject or food between stages. These bare clauses permit contextually
understood participants. The morphology is \ird{pjašt-é} (eat-\Proh{}),
\ird{pjašt-ku} (eat-\Neg{}\Hort{}), and
\ird{pjašt-ir} (eat-\Neg{}\Irr{}). Let $D$, $H$, and $R$ stand for
imperative, hortative, and irrealis force respectively. The evaluation stages
belong to the course of action under direction/proposal, or to the relevant
hypothetical situations under irrealis; these clauses do not assert an actual
prior episode of eating.

\begin{table}[ht]
\footnotesize
\caption{Negative Mood diagnostics: fixed action alternatives.}
\label{tab:th-negative-mood-diagnostics}
\begin{tabularx}{\textwidth}{L l l}
\toprule
\textsc{surface string} & \textsc{ordinary scope} & \textsc{status} \\
\midrule
\ird{Oh vy pjašté.} & $D(A(V(\neg P)))$ & M \\
\ird{Oh so pjašté.} & $D(A(S(\neg P)))$ & L \\
\ird{Kamo vy pjašté.} & $D(E(V(\neg P)))$ & M \\
\ird{Kamo so pjašté.} & $D(E(S(\neg P)))$ & L \\
\ird{Oh vy pjaštku.} & $H(A(V(\neg P)))$ & M \\
\ird{Oh so pjaštku.} & $H(A(S(\neg P)))$ & L \\
\ird{Kamo vy pjaštku.} & $H(E(V(\neg P)))$ & M \\
\ird{Kamo so pjaštku.} & $H(E(S(\neg P)))$ & L \\
\ird{Oh vy pjaštir.} & $R(A(V(\neg P)))$ & L \\
\ird{Oh so pjaštir.} & $R(A(S(\neg P)))$ & L \\
\ird{Kamo vy pjaštir.} & $R(E(V(\neg P)))$ & L \\
\ird{Kamo so pjaštir.} & $R(E(S(\neg P)))$ & L \\
\bottomrule
\end{tabularx}
\end{table}

The \ird{vy} rows urge, propose, or hypothetically describe cessation of
eating; the \ird{so} rows continued non-eating. M means licensed but
pragmatically marked: a cessation directive or proposal needs an established
course to stop, rather than a neutral preventative setting. It becomes
appropriate with an instruction such as \trsl{From now on, stop eating too}
for additive \ird{vy}, or \trsl{Only stop eating now; continue the other
activities} for high exclusive \ird{vy}. For \ird{so}, corresponding contexts
are \trsl{Continue refraining from eating too} and \trsl{Only continue
refraining from eating}; the additive or exclusive comparison is with the same
other actions. All twelve rows inherit the stage and alternative requirements
above, interpreted inside $D/H/R$. The lower exclusive option, when its
stagewise background is satisfied, is $M(V(E(\neg P)))$ or
$M(S(E(\neg P)))$ for $M\in\{D,H,R\}$.

Excluded analyses are $\neg M(C(\mathrm{PHAS}(P)))$ and
$M(C(\neg\mathrm{PHAS}(P)))$ as readings of these negative Mood patterns:
they place the single negation outside its licensed input position.
\ird{*Oh so zapjašté}, \ird{*Oh so zapjaštku}, and
\ird{*Oh so zapjaštir} are excluded attempts to add indicative negation to
negative Mood. Lowering additive \ird{oh} to derive
$M(\mathrm{PHAS}(A(\neg P)))$ is likewise excluded by the present rules.
No change to the established surface order, aspect quantity, or negative
suffix paradigms follows from these diagnostics.
